\section{Limitations and Non-Claims}
\label{sec:limitations}

The scope of the evidence is intentionally narrow.

\begin{itemize}
  \item We probe a pre-existing prompting channel. We do not fine-tune student
  models here, and the results do not establish what is necessary or sufficient
  for training-time subliminal learning.

  \item The layerwise output-head readout is correlational. The residual-state
  patch establishes sufficiency of a full assistant-position state, but it does
  not identify a unique causal direction, circuit, neuron set, attention head,
  or algorithm. Nor does it establish necessity.

  \item One same-release Llama 8B--70B pair does not establish a universal
  scaling law. The 1B/3B/8B descriptive ladder mixes Llama 3.2 and 3.1 releases
  and is not used as the clean causal contrast.

  \item Two small Qwen models do not establish a tokenizer-family law. A larger
  crossed design is required to separate tokenizer, architecture, training
  corpus, and scale.

  \item The final animal-specific contrast reaches correlation one algebraically
  after removing the shared softmax denominator. Its interpretation concerns
  the pre-final trajectory and AUC, not the endpoint.

  \item Width-1 Qwen contains only ten strings and cannot support a powered
  inference; it is reported descriptively. Width-specific heterogeneity also
  warns against extrapolating the width-3 primary to arbitrary sequences.

  \item Bootstrap intervals summarize variation across the fixed 18 animals.
  The causal bootstrap additionally resamples the fixed unordered number-pair
  clusters. Neither implies random sampling from every semantic concept or model
  family.

  \item The novelty statement is based on a documented literature search. It is
  not an absolute claim that no concurrent or unindexed work has run a similar
  experiment.
\end{itemize}

Within those boundaries, the core result is stable: static output geometry,
tokenizer atomicity, contextual answer readability, and causal state
transmission are empirically distinct measurements and should not be treated as
one scalar mechanism.
